\chapter{Experimental Methodology}
\label{ch:experimental_methodology}

This chapter details the experimental design, benchmarking framework, and statistical methods used to evaluate the search performance and playing strength of Kronos.

\section{Automated Tournament Scheduler}
The benchmarking suite runs tournaments using a Node.js script managed by our \texttt{TournamentRunner}. The scheduler supports paired opening match execution. For each match, the engine variants play two games using a shared opening FEN: once with Engine A playing White, and once with Engine A playing Black. This paired arrangement isolates playing strength evaluations from opening book advantages (such as the standard first-move advantage of White). 

To accelerate tournament execution, the framework supports parallel worker threads. However, this concurrency represents a systems-level trade-off. Although the host hardware has 12 logical processors, each V8 engine process exhibits significant untracked memory growth outside the managed heap. The measured peak RSS for a single Depth~6 tournament worker is 1,713\,MB (Table~\ref{tab:memory_profile}), despite V8 reporting only 445\,MB of active heap usage---a 3.5$\times$ ratio. This gap is attributable to V8 JIT-compiled code segments, external array buffers, \texttt{chess.js} internal string and object allocations that escape V8 heap accounting, and Node.js worker thread overhead. Critically, this untracked RSS growth is cumulative over tournament execution (hundreds of games) and does not release predictably.

While the nominal aggregate RSS of three workers (approximately 5.1\,GB) appears well within the 16\,GB physical memory budget, empirical testing revealed that launching a fourth concurrent worker triggered V8 heap paging and GC sweep distortion. We attribute this to the cumulative untracked memory growth compounding across workers: at four workers, the combined untracked allocations (JIT caches, external buffers, and OS-level per-thread overhead) exceeded available physical memory after sustained tournament execution, despite V8's managed heap remaining within bounds. Capping concurrency at three workers (referred to as Cap~3) was the empirically determined threshold that prevented measurement distortion across full tournament runs.


\begin{algorithm2e}[tp]
\caption{Tournament Scheduler with SPRT}
\label{alg:tournament_scheduler}
\DontPrintSemicolon
\Input{Engine Config A, Engine Config B, Opening FENs list $F$, SPRT bounds $[\eta_0, \eta_1]$, Max Games $M$}
\Output{Trinomial Elo rating delta, SPRT Decision}
\BlankLine
$wins \leftarrow 0$, $losses \leftarrow 0$, $draws \leftarrow 0$, $N \leftarrow 0$\;
$checkpoint \leftarrow \text{LoadCheckpoint()}$\;
\If{$checkpoint$ exists}{
    $wins, losses, draws, N \leftarrow \text{RestoreState}(checkpoint)$\;
}
\While{$N < M$}{
    $fen \leftarrow F[N / 2]$\;
    \If{$N \pmod 2 == 0$}{
        $result \leftarrow \text{PlayGame}(\text{Config A as White}, \text{Config B as Black}, fen)$\;
    }
    \Else{
        $result \leftarrow \text{PlayGame}(\text{Config B as White}, \text{Config A as Black}, fen)$\;
        $result \leftarrow \text{Flip}(result)$ \tcp*{Invert outcome for Config A perspective}
    }
    $wins, losses, draws \leftarrow \text{UpdateCounts}(wins, losses, draws, result)$\;
    $N \leftarrow N + 1$\;
    $\text{SaveCheckpoint}(wins, losses, draws, N)$\;
    $LLR \leftarrow \text{CalculateLLR}(wins, losses, draws)$\;
    \If{$LLR \ge \eta_1$}{
        \Return{\text{ComputeElo}(wins, losses, draws), $\text{Accepted } (H_1)$}\;
    }
    \If{$LLR \le \eta_0$}{
        \Return{\text{ComputeElo}(wins, losses, draws), $\text{Rejected } (H_1)$}\;
    }
}
\Return{\text{ComputeElo}(wins, losses, draws), \text{Inconclusive}}\;
\end{algorithm2e}


The scheduler shuffles the opening FEN library (\texttt{openings.json}) using a custom Linear Congruential Generator (LCG) Pseudo-Random Number Generator (PRNG) initialized with a fixed seed. This ensures that different tournament runs can be reproduced under identical game configurations. The execution pipeline runs from raw configuration specs to final web-based dashboards, as detailed in Figure~\ref{fig:benchmark_framework}.

\begin{figure}[htbp]
\centering
\begin{tikzpicture}[
    node distance=1.0cm and 2.0cm,
    block/.style={draw=blue!60!black, rectangle, minimum height=2.8em, minimum width=10em, align=center, rounded corners=3pt, fill=blue!10, font=\small\sffamily},
    arrow/.style={draw, -latex, thick, gray!80!black}
]
    \node [block] (source) {Position Suite};
    \node [block, right=of source] (runner) {Pipeline Manager};
    \node [block, right=of runner] (ordo) {Rating Evaluator};
    \node [block, below=of runner, yshift=-0.2cm] (output) {Report Generator};

    \draw [arrow] (source) -- (runner);
    \draw [arrow] (runner) -- (ordo);
    \draw [arrow] (ordo.south) |- (output.east);
    \draw [arrow] (runner) -- (output);
\end{tikzpicture}
\caption[Kronos Benchmarking Framework Data Flow]{Kronos Benchmarking Framework Data Flow.}
\label{fig:benchmark_framework}
\end{figure}

% Source Material: Summarized from benchmark/docs/README.md


\section{Execution Integrity and Verification}
To ensure that engine matches reflect correct chess logic, the runner monitors game play. We verified execution integrity by ensuring 100\% legal moves under FIDE rules, strict color alternation (W $\to$ B $\to$ W), and correct draw trigger validation for stalemates, threefold repetitions, and the 50-move rule. The framework logs game records to standard PGN files, and we confirmed that all exported PGNs parse and replay cleanly in standard chess readers. The data flow of the validation pipeline is illustrated in Figure~\ref{fig:validation_pipeline}.

\begin{figure}[htbp]
\centering
\begin{tikzpicture}[
    node distance=1.0cm and 2.0cm,
    block/.style={draw=blue!60!black, rectangle, minimum height=2.8em, minimum width=10em, align=center, rounded corners=3pt, fill=blue!10, font=\small\sffamily},
    arrow/.style={draw, -latex, thick, gray!80!black}
]
    \node [block] (games) {Game Results (W/D/L)};
    \node [block, right=of games] (stats) {SPRT Calculator};
    \node [block, right=of stats] (val) {Validation Report};

    \draw [arrow] (games) -- (stats);
    \draw [arrow] (stats) -- (val);
\end{tikzpicture}
\caption[Kronos Validation Pipeline]{Kronos Validation Pipeline.}
\label{fig:validation_pipeline}
\end{figure}

% Source Material: Summarized from STATISTICAL_VALIDATION_REPORT.md


\section{Statistical Heuristics and Formulas}
We compute tournament ratings and confidence intervals using the Bradley-Terry logistic model. The expected score $E_A$ (representing the probability of player A defeating player B, plus half the probability of a draw) and the rating updates are computed using standard formulas:

% Elo rating update formula and expected score formula
\begin{equation}
    E_A = \frac{1}{1 + 10^{(R_B - R_A)/400}}
\end{equation}
\begin{equation}
    R'_A = R_A + K(S_A - E_A)
\end{equation}
% Where:
% R_A, R_B are current ratings of player A and B.
% E_A is the expected score of player A.
% R'_A is the updated rating of player A.
% S_A is the actual score (1 for win, 0.5 for draw, 0 for loss).
% K is the rating development factor.


To evaluate rating differences and confidence intervals from match scores, we compute the trinomial sample variance and standard error:

% Statistical Confidence Intervals for Elo ratings
\begin{equation}
    \sigma_{score} = \sqrt{\frac{p_{win} + p_{loss} - (p_{win} - p_{loss})^2}{N}}
\end{equation}
\begin{equation}
    CI_{95\%} = \hat{S} \pm 1.96 \cdot \frac{\sigma_{score}}{\sqrt{N}}
\end{equation}
where $p_{win}$ and $p_{loss}$ are the sample probabilities of winning and losing, $N$ is the total number of games played, and $\hat{S}$ is the sample mean expected score defined as $\hat{S}=p_{win}+0.5\cdot p_{draw}$.


We apply Wald's Sequential Probability Ratio Test (SPRT) to determine whether a search enhancement provides a statistically significant playing strength increase.

The SPRT routine evaluates the log-likelihood ratio (LLR) of match outcomes against a null hypothesis $H_0$ ($\text{Elo} \le 0$) and an alternative hypothesis $H_1$ ($\text{Elo} \ge 10$). The runner computes the LLR after each game pair. If the LLR exceeds the upper boundary ($\eta_1 = \ln \frac{1-\beta}{\alpha}$), the tournament terminates and accepts $H_1$. Conversely, if the LLR falls below the lower boundary ($\eta_0 = \ln \frac{\beta}{1-\alpha}$), the scheduler terminates and accepts $H_0$. Under target error rates of $\alpha = 0.05$ and $\beta = 0.05$, these boundary thresholds correspond to $[-2.944, 2.944]$. Here, $\ln$ represents the natural logarithm. Frequentist $p$-values are not computed under this sequential testing framework; instead, the hypothesis test terminates strictly when the LLR crosses these Wald boundaries.


When running fixed-length calibration tournaments rather than SPRT matchups, the scheduler evaluates Elo boundaries using trinomial distribution updates. In this mode, execution halts early if the 95\% confidence radius drops below $\pm 60$ Elo or stabilizes within a range of $5$ Elo over a 10-game window.


To quantify the search space compression achieved by each search heuristic, we define the \textbf{Node Reduction (\%)} metric:
\begin{equation}
    \text{Node Reduction (\%)} = \frac{N_{\text{disabled}} - N_{\text{baseline}}}{N_{\text{disabled}}} \times 100\%
\end{equation}
where $N_{\text{disabled}}$ is the number of visited nodes with the heuristic deactivated, and $N_{\text{baseline}}$ is the number of visited nodes with the heuristic active (Baseline engine). This represents the percentage of the search tree pruned specifically by activating that heuristic.

Finally, the Effective Branching Factor (EBF) measures search tree expansion:

% Effective Branching Factor (EBF) equation
\begin{equation}
    N(d) = \sum_{i=0}^{d} (b^*)^i = \frac{(b^*)^{d+1} - 1}{b^* - 1}
\end{equation}
\begin{equation}
    b^* \approx \frac{N_d}{N_{d-1}}
\end{equation}
% Where:
% N(d) is the total nodes evaluated at depth d.
% b^* is the effective branching factor.



\subsection{Metric Definitions}
To ensure clarity and consistency in our evaluation, we define the primary quantitative metrics used in this work as follows:
\begin{itemize}
    \item \textbf{Internal Bradley-Terry Self-Play Rating (Elo):} The relative playing strength calculated within the closed tournament pool of Kronos engine variants, fitted using the Bradley-Terry logistic model and anchored at 1,000 Elo for the Baseline Minimax engine.
    \item \textbf{Relative Gain (Elo):} The incremental change in the internal self-play rating between adjacent search depths ($d$ vs.\ $d-1$ plies) in the closed self-play pool.
    \item \textbf{Node Reduction (\%):} The percentage of search tree size compression achieved by enabling a search heuristic, defined as $(N_{\text{disabled}} - N_{\text{baseline}}) / N_{\text{disabled}} \times 100\%$, where $N_{\text{disabled}}$ is the number of visited nodes with the heuristic deactivated, and $N_{\text{baseline}}$ is the number of visited nodes with the heuristic active.
    \item \textbf{Speedup:} The ratio of search execution speed (based on time elapsed per search) of a disabled variant compared to the active baseline variant ($T_{\text{disabled}} / T_{\text{baseline}}$).
    \item \textbf{Effective Branching Factor (EBF):} The geometric mean number of branches expanded per ply, calculated as $b = N^{1/d}$ for search nodes $N$ at depth $d$.
    \item \textbf{Relative Depth-Limited Matchup Differential (Elo):} The measured playing strength difference obtained by benchmarking the final Kronos engine (Depth 6) against fixed-depth reference configurations of Stockfish 18. These differentials reflect performance against crippled Stockfish configurations and should not be mapped to FIDE or CCRL rating scales.
\end{itemize}

\section{Checkpointing and Memory Profiling}
Tournaments run for hundreds of games, taking hours to execute. The framework implements two mechanisms to ensure robustness:
\begin{itemize}
    \item \textbf{JSON Checkpointing:} The tournament state is saved to a checkpoint JSON file every 10 games. The checkpoint records the game index, win/loss/draw counts, active PRNG seed index, color-flip flags, and telemetry logs. If a crash or interruption occurs, the runner reads the JSON file and resumes execution from the last saved state without loss of progress.
    \item \textbf{Memory Profiler:} The scheduler monitors V8 heap allocations and transposition table footprints every 5 games. The profiler logs system memory usage (RSS, heap total, heap used) to \texttt{MEMORY\_PROFILE.md} during execution. This tracking monitors memory leaks that could degrade search performance during long runs.
\end{itemize}

To specify the configurations tested, Table~\ref{tab:experimental_configs} shows the parameter flags defined in the engine config files.

\begin{table}[htbp]
\centering
\caption[Experimental Configuration Parameters]{Experimental configuration parameters for the benchmark run profiles.}
\label{tab:experimental_configs}
\begin{tabularx}{\textwidth}{lXc}
\toprule
\textbf{Parameter Flag} & \textbf{Search Module Influence} & \textbf{Default} \\
\midrule
\texttt{useAlphaBeta} & Toggles baseline minimax vs.\ alpha-beta pruning & True \\
\texttt{useIterativeDeepening} & Controls incremental deepening search loops & True \\
\texttt{useMoveOrdering} & Enables MVV-LVA, killer moves, and history heuristics & True \\
\texttt{useTranspositionTable} & Toggles Zobrist-keyed transposition cache & True \\
\texttt{useQuiescence} & Restricts leaf search to quiet positions & True \\
\texttt{usePieceSquareTables} & Controls positional evaluation mapping & True \\
\bottomrule
\end{tabularx}
\end{table}

